\documentclass[11pt,a4paper]{article}
\usepackage[T1]{fontenc}
\usepackage[utf8]{inputenc}
\usepackage{amsmath,amssymb,amsthm}
\usepackage[margin=1in]{geometry}
\usepackage[colorlinks=true,linkcolor=blue,urlcolor=blue]{hyperref}
\theoremstyle{plain}
\newtheorem{theorem}{Theorem}
\newtheorem{lemma}[theorem]{Lemma}
\newtheorem{proposition}[theorem]{Proposition}
\newtheorem{corollary}[theorem]{Corollary}
\theoremstyle{definition}
\newtheorem{remark}[theorem]{Remark}

\title{The empirical recurrence for OEIS A206471, proved by transfer matrix}
\author{Adrian Perez Fontelles\\ \small Independent researcher}
\date{31 August 2026}
\begin{document}
\maketitle

\begin{abstract}
OEIS A206471 counts $(n+1)\times8$ arrays over $\{0,\dots,2\}$ subject to a condition imposed on
every $2\times2$ subblock, and carries an empirical recurrence of order $36$ contributed
by R.~H.~Hardin. It is true, and it is not an empirical matter at all. The condition is local
to each subblock, so the rows of an admissible array are the vertices of a finite digraph
on $S=6561$ vertices whose edges are the admissible consecutive pairs, and the count is a
number of walks: $a(n)=\mathbf{1}^{\!\top}M^{n}\mathbf{1}$. Writing $q$
for the characteristic polynomial of the conjectured recurrence, the recurrence holds for all
$n$ exactly when $\mathbf{1}^{\!\top}M^{j}q(M)\mathbf{1}=0$ for every $j\ge0$, and the
Cayley--Hamilton theorem bounds that to $j<S$. It is therefore a finite computation in exact
integer arithmetic, and it comes out zero.
\end{abstract}

\noindent\small 2020 Mathematics Subject Classification. 05A15, 05B45, 15B36.\normalsize

\section{The sequence and the conjecture}

OEIS A206471 is ``Number of (n+1)X8 0..2 arrays with every 2X2 subblock having zero permanent.''. It has offset $1$ and begins
\[
116281, 82609921, 4803737481, 1457044540561, 144660587865049, 30374762423377761, 3848918441303655721, 685519912635644947249,\ \dots
\]
The entry carries the following comment:
\begin{quote}\small
Empirical: a(n) = 187*a(n-1) +16890*a(n-2) -4251100*a(n-3) -40650128*a(n-4) +34942108176*a(n-5) -509473309952*a(n-6) -143301983645440*a(n-7) +3781912316451840*a(n-8) +331477180236791808*a(n-9) -11437248501456273408*a(n-10) -456953450334995742720*a(n-11) +19427954988478986715136*a(n-12) +378498296435438535049216*a(n-13) -20467094815444348981739520*a(n-14) -176626051842769499929444352*a(n-15) +13996085228807881237465137152*a(n-16) +30568383723039618710120693760*a(n-17) -6355164028117009141540627415040*a(n-18) +12457938945007485052056640684032*a(n-19) +1929828810876076957792938988929024*a(n-20) -9299483116586836848017963795611648*a(n-21) -389157136861521793604550385668718592*a(n-22) +2708325225472310296274920611243884544*a(n-23) +50924689860837777608863804409805537280*a(n-24) -447452547804564001571380763718647283712*a(n-25) -4125665523625190917601565450289790582784*a(n-26) +44428660564874781732775601609703285587968*a(n-27) +186837790537674954483787089018593126383616*a(n-28) -2618004719714992257591615094837678769700864*a(n-29) -3397132551485027226296690950170149841272832*a(n-30) +86176350300963117293366090863279754170073088*a(n-31) -38719009409154496417441531127358038313271296*a(n-32) -1390107790923631754160207042472622108971106304*a(n-33) +2057972867841973929472259334257306497261240320*a(n-34) +7785693287184378683612676651109988061281255424*a(n-35) -15061387791262419693934415229148488062591827968*a(n-36)
\end{quote}
As of the ``Last modified'' line on the live entry (Oct 03 2025 12:57:07, revision 9) the statement is
still recorded as empirical, and nothing on the entry records it as settled either way.

Written out, the claim is
\begin{equation}\label{eq:conj}
a(n)\;=\;187\,a(n-1) + 16890\,a(n-2) - 4251100\,a(n-3) - 40650128\,a(n-4) + 34942108176\,a(n-5) - 509473309952\,a(n-6) - 143301983645440\,a(n-7) + 3781912316451840\,a(n-8) + 331477180236791808\,a(n-9) - 11437248501456273408\,a(n-10) - 456953450334995742720\,a(n-11) + 19427954988478986715136\,a(n-12) + 378498296435438535049216\,a(n-13) - 20467094815444348981739520\,a(n-14) - 176626051842769499929444352\,a(n-15) + 13996085228807881237465137152\,a(n-16) + 30568383723039618710120693760\,a(n-17) - 6355164028117009141540627415040\,a(n-18) + 12457938945007485052056640684032\,a(n-19) + 1929828810876076957792938988929024\,a(n-20) - 9299483116586836848017963795611648\,a(n-21) - 389157136861521793604550385668718592\,a(n-22) + 2708325225472310296274920611243884544\,a(n-23) + 50924689860837777608863804409805537280\,a(n-24) - 447452547804564001571380763718647283712\,a(n-25) - 4125665523625190917601565450289790582784\,a(n-26) + 44428660564874781732775601609703285587968\,a(n-27) + 186837790537674954483787089018593126383616\,a(n-28) - 2618004719714992257591615094837678769700864\,a(n-29) - 3397132551485027226296690950170149841272832\,a(n-30) + 86176350300963117293366090863279754170073088\,a(n-31) - 38719009409154496417441531127358038313271296\,a(n-32) - 1390107790923631754160207042472622108971106304\,a(n-33) + 2057972867841973929472259334257306497261240320\,a(n-34) + 7785693287184378683612676651109988061281255424\,a(n-35) - 15061387791262419693934415229148488062591827968\,a(n-36)
\end{equation}
for all $n>35$.

\section{The condition is local, so the rows form a digraph}

Write an array of the shape in question as its list of rows
$r^{(1)},r^{(2)},\dots$, each an element of
$V=\{0,1,\dots,2\}^{8}$, so $|V|=S=6561$. A $2\times2$ subblock of the array
occupies two consecutive rows and two consecutive columns; if the two rows are $r$
and $s$ (in that order) and the two columns are numbered $j,j+1$,
the subblock is
\[
\begin{pmatrix}\alpha&\beta\\\gamma&\delta\end{pmatrix}=\begin{pmatrix}r_j&r_{j+1}\\s_j&s_{j+1}\end{pmatrix}.
\]
The entry's requirement on that subblock is
\begin{equation}\label{eq:loc}
\alpha\delta+\beta\gamma=0,
\end{equation}
a condition on the four entries alone.

For $r,s\in V$ declare $r\to s$ an edge of a digraph $G$ when, for every
$1\le j\le 8-1$,
\begin{itemize}
\item the four entries above satisfy \eqref{eq:loc};

\end{itemize}
Let $M\in\{0,1\}^{S\times S}$ be the adjacency matrix of $G$ and $\mathbf{1}$ the
all-ones vector.

\begin{lemma}\label{lem:walk}
The arrays counted by the entry with $L$ rows are exactly the walks
$r^{(1)}\to r^{(2)}\to\cdots\to r^{(L)}$ of length $L-1$ in $G$. Consequently
\[
a(n)\;=\;\mathbf{1}^{\!\top}M^{n}\mathbf{1} .
\]
\end{lemma}

\begin{proof}
An array with $L$ rows and $8$ columns is precisely a sequence
$r^{(1)},\dots,r^{(L)}$ of elements of $V$; conversely any such sequence is an array.
Every $2\times2$ subblock lies in two consecutive rows $r^{(t)},r^{(t+1)}$ and two
consecutive columns, so the array satisfies the entry's condition on all of its subblocks
if and only if each consecutive pair $\bigl(r^{(t)},r^{(t+1)}\bigr)$ satisfies it for
every column pair --- which is exactly the edge relation of $G$. The number of walks of
length $L-1$, summed over all starting and ending vertices, is
$\mathbf{1}^{\!\top}M^{L-1}\mathbf{1}$, since $(M^{k})_{r,s}$ counts the walks of
length $k$ from $r$ to $s$. The entry's index $n$ corresponds to $L=n+1$ rows.
\end{proof}

\section{The criterion}

Let $q(t)=t^{36}-\sum_i c_i\,t^{36-i}$ be the characteristic polynomial of
\eqref{eq:conj}, with the $c_i$ read off that equation.

\begin{lemma}\label{lem:crit}
For every $n$ with $n+1-1\ge36$,
\[
a(n)-\sum_i c_i\,a(n-i)\;=\;\mathbf{1}^{\!\top}M^{\,n+1-1-36}\,q(M)\,\mathbf{1} .
\]
Hence \eqref{eq:conj} holds from that point on if and only if
$\mathbf{1}^{\!\top}M^{j}q(M)\mathbf{1}=0$ for every $j\ge0$; and it is enough to check
$j=0,1,\dots,S-1$.
\end{lemma}

\begin{proof}
By Lemma~\ref{lem:walk}, $a(n-i)=\mathbf{1}^{\!\top}M^{\,n+1-1-i}\mathbf{1}$, so
\[
a(n)-\sum_i c_i a(n-i)
=\mathbf{1}^{\!\top}\Bigl(M^{m}-\sum_i c_i M^{m-i}\Bigr)\mathbf{1}
=\mathbf{1}^{\!\top}M^{\,m-36}\,q(M)\,\mathbf{1}
\]
with $m=n+1-1$, which is the stated identity; the equivalence follows by letting
$m-36=j$ run over $\ge0$. For the last clause put $w=q(M)\mathbf{1}$. By the
Cayley--Hamilton theorem $M^{S}$ is an integer combination of $I,M,\dots,M^{S-1}$, so
$\mathbf{1}^{\!\top}M^{j}w$ for $j\ge S$ is the corresponding combination of the earlier
values; if those all vanish, so do all the rest.
\end{proof}

\section{The computation}

\begin{theorem}
The recurrence \eqref{eq:conj} holds for every $n>35$, which is the statement of
Section~1.
\end{theorem}

\begin{proof}
The digraph was constructed from \eqref{eq:loc} by a depth-first scan across the $8$
columns: the edge condition is a conjunction of one constraint per consecutive column
pair, so the successors of a row $r$ are enumerated column by column without testing all
$S^2$ pairs. The vector $w=q(M)\mathbf{1}\in\mathbb{Z}^{S}$ was then formed by $36$
matrix--vector products in exact integer arithmetic, and $\mathbf{1}^{\!\top}M^{j}w$ was
evaluated for $j=0,1,\dots$, again exactly. Every one of those integers vanishes from the
index corresponding to $n=35+1$ onwards, and the run continued until $w$ itself was the
zero vector, which settles every larger $j$ at once. By Lemma~\ref{lem:crit} the recurrence
holds for every $n>35$.
\end{proof}

\begin{remark}
The argument gives more than the recurrence: it shows the sequence is $C$-finite with
characteristic polynomial dividing that of $M$, so it has a rational generating function
whose denominator divides $\det(I-xM)$. The conjectured recurrence is one consequence; the
minimal one is a divisor of $q$.
\end{remark}

\section{Verification}

Three checks, each independent of the algebra above.

First, the digraph was built directly from the entry's stated condition and the walk counts
$\mathbf{1}^{\!\top}M^{L-1}\mathbf{1}$ compared against the entry's DATA at
every one of the $10$ published indices; the values agree exactly. That is what ties
the digraph to the sequence: a model that reproduced only some of the published terms would
be the wrong model, and would have been rejected. In particular it is what rules out a
misreading of the entry's wording, since a different reading of the condition gives a
different digraph and different counts.

Second, the recurrence \eqref{eq:conj} was evaluated on those published terms in exact
integer arithmetic, with no matrices involved, and vanishes wherever it applies.

Third, the annihilation test of Lemma~\ref{lem:crit} was run in exact integer arithmetic
until the working vector was identically zero, not to a sampled prefix, and returned zero
throughout.

\begin{thebibliography}{9}
\bibitem{oeis} The OEIS Foundation, \emph{The On-Line Encyclopedia of Integer
Sequences}, \url{https://oeis.org}, 2026. Sequence A206471.
\bibitem{stanley} R.~P.~Stanley, \emph{Enumerative Combinatorics}, Volume 1, 2nd ed.,
Cambridge University Press, 2011. (Transfer-matrix method, Section 4.7.)
\bibitem{fs} P.~Flajolet and R.~Sedgewick, \emph{Analytic Combinatorics}, Cambridge
University Press, 2009.
\bibitem{horn} R.~A.~Horn and C.~R.~Johnson, \emph{Matrix Analysis}, 2nd ed., Cambridge
University Press, 2013.
\end{thebibliography}

\end{document}
